%%%%%%%%%%%%%%%%%%%%%%%%%%%%%%%%%%%%%%%%%%%%%%%%%%%%%%%%%%%%%%%%%%%%%%
% Lớp 9 | Chương I. Phương trình và hệ hai phương trình bậc nhất hai ẩn | Giải bài toán bằng cách lập hệ phương trình
% Dạng 9D1?3-2: Toán năng suất, làm chung – làm riêng
%%%%%%%%%%%%%%%%%%%%%%%%%%%%%%%%%%%%%%%%%%%%%%%%%%%%%%%%%%%%%%%%%%%%%%


% nguồn: 2026-11-15_09_de_tham_khao_quan_3_de_1_va_de_2 (nhập 15/11/2026)
\begin{bt}%[9D1V3-2]%[ma:b3d00470]
% nguon: Đề tham khảo Tuyển sinh 10 TP.HCM 2025-2026 – Quận 3 – Đề 2 – Bài 6
Nhân dịp ngày Nhà giáo Việt Nam, nhà trường tổ chức cuộc thi kéo co. Giáo viên chủ nhiệm lớp 9B đã chọn ra $\dfrac{1}{4}$ số học sinh nam và $\dfrac{1}{3}$ số học sinh nữ của lớp tham gia kéo co, biết rằng số học sinh nam và số học sinh nữ tham gia kéo co là bằng nhau. Toàn bộ $20$ học sinh còn lại của lớp tham gia đội cổ động. Hãy tìm số học sinh nam và số học sinh nữ của lớp 9B.
\loigiai{
Gọi số học sinh nam và số học sinh nữ của lớp 9B lần lượt là $x$ và $y$ ($x, y \in \mathbb{N}^*$, $x$ chia hết cho $4$ và $y$ chia hết cho $3$).\\
Số học sinh nam tham gia kéo co là $\dfrac{1}{4}x$ (học sinh).\\
Số học sinh nữ tham gia kéo co là $\dfrac{1}{3}y$ (học sinh).\\
Vì số học sinh nam và số học sinh nữ tham gia kéo co bằng nhau nên ta có phương trình:
$\dfrac{1}{4}x = \dfrac{1}{3}y \Leftrightarrow 3x - 4y = 0$.\\
Số học sinh nam làm cổ động viên là: $x - \dfrac{1}{4}x = \dfrac{3}{4}x$ (học sinh).\\
Số học sinh nữ làm cổ động viên là: $y - \dfrac{1}{3}y = \dfrac{2}{3}y$ (học sinh).\\
Vì đội cổ động có tổng cộng $20$ học sinh nên ta có phương trình:
$\dfrac{3}{4}x + \dfrac{2}{3}y = 20 \Leftrightarrow 9x + 8y = 240$.\\
Ta có hệ phương trình:
$\begin{cases} 3x - 4y = 0 \\ 9x + 8y = 240 \end{cases} \Leftrightarrow \begin{cases} 6x - 8y = 0 \\ 9x + 8y = 240 \end{cases} \Leftrightarrow \begin{cases} 15x = 240 \\ y = \dfrac{3}{4}x \end{cases} \Leftrightarrow \begin{cases} x = 16 \\ y = 12 \end{cases}$ (thoả mãn điều kiện).\\
Vậy lớp 9B có $16$ học sinh nam và $12$ học sinh nữ.
}
\end{bt}
